\chapter{引言}

\section{实验背景及意义}

哈希表通过哈希函数把关键字映射到数组位置\cite{bender2024sosa}，常用于数据库索引、键值存储、编译器符号表和运行时系统\cite{chang2006bigtable,dong2021rocksdb}。在这类场景中，查询路径通常较短，工程实现也相对成熟。但当数据规模扩大后，哈希表为维持查询和更新所保存的辅助信息会逐渐占据不可忽略的空间。系统除存放键或键的部分信息外，还要记录冲突处理、删除标记、元素重定位和键恢复所需的元数据。面向键值存储的哈希索引优化也是国内已有工作的关注方向之一，相关研究从键摘要压缩、热点感知调度等角度改进空间利用和探查路径\cite{wang2024hotaware}。因此，评价哈希表时需要同时考察平均查询时间和每个键背后的元数据成本。

相较于链地址法，开放寻址哈希表把元素直接放在连续数组中，查询时不需要沿指针跳转，也更容易利用缓存局部性。它的代价也比较明确：负载因子升高后，冲突和探测长度会增加，工程实现往往需要预留空槽来降低冲突概率。删除操作还会带来额外限制。删除后的空槽如果直接释放，可能破坏后续查询路径；若采用回填或重定位，又必须维护被移动元素的状态。因此本文关注的核心问题是：在查询、插入和删除时间可控的前提下，尽量降低元素定位、更新和键恢复所需的辅助存储，并观察相关理论结构落到原型系统后产生的维护成本。

Bender 等人的论文\cite{bender2022otsh} 从理论层面给出了更详细的描述：若插入和删除允许承担 \(O(k)\) 次元素交换代价，即一次更新过程中最多搬移 \(O(k)\) 个元素，则每个元素的额外空间可以从 \(O(\log\log n)\) 比特降至 \(O(\log^{(k)}n)\) 比特，同时查询时间仍保持 \(O(1)\)，这里 \(\log^{(k)}n\) 表示对 \(n\) 连续取 \(k\) 次对数。该结果表明，开放寻址哈希表的空间开销并不只由负载率决定，还与更新过程中的元素搬移次数以及从表中恢复原始键所需的编码信息有关。本文基于这一结构完成原型实现与实验分析，重点观察上述理论权衡在工程实现中会转化为哪些具体开销。

\section{国内外研究现状}

\subsection{经典冲突处理}

常见哈希表实现大体上能分成链地址法和开放寻址法\cite{bender2024sosa}两类，链地址法一般借助链表或桶来处理冲突，写法比较直接，插入删除也相对好理解，可每个元素常要额外存指针，内存怎么分配也会打乱连续访问的效果；开放寻址法则把元素放进数组槽里，访问时数据更集中，但也更容易受负载率、探针顺序以及删除策略的影响，线性探测和它的几种改法通过连续探查换取更好的缓存表现\cite{pagh2004linear,thorup2012tabulation}，Robin Hood 哈希这类方法主要想让不同元素的查找距离尽量接近\cite{conway2018}，这些做法一到删除多或装填偏高的情形就麻烦，用墓碑占位会把查询拖长，若马上回填又得保证被挪走的元素仍能按原路径查到。

\subsection{踢出式与完美哈希}

布谷鸟哈希\cite{pagh2004cuckoo,lehmann2023sichash}会给每个元素准备多个候选槽，冲突时通过踢出把查找变短，但负载升高后仍要处理踢出链和整体重哈希这类问题；动态完美哈希\cite{dietzfelbinger1994}在最坏情况下也能把查询时间压在常数级，给了经典的上界结果，最小完美哈希\cite{lehmann2026mphsurvey,lehmann2024shockhash,hermann2024phobic,pibiri2021pthash,esposito2020recsplit}主要服务静态集合，查找保证更强，更适合关键词不再变、查多改少的场合，其后又有工作去压建表开销、压缩编码、缩短找种子的过程\cite{lehmann2025combined}，这类结构在不动态变化的数据上省空间，可一旦放进需要频繁增删的键值存储，还得另加办法才能做增量更新，重构失败时也要另行处理。

\subsection{紧凑存储与空间--时间权衡}

只看负载率还说不清紧凑哈希表到底费多少空间，因为系统还得留着辅助信息来支撑查询、更新和把键找回来，Bloom 过滤器\cite{bloom1970,broder2003bloom}这类近似判断结构，加上动态检索和过滤器下界方面的工作\cite{kuszmaul2024filters}，都说明紧凑结构一旦要支持更新，状态一变就要多占信息；紧凑检索结构\cite{pagh2008retrieval}则演示了怎样用较小空间存和键有关的值，Bender 等人提出的时间/空间权衡框架\cite{bender2022otsh}把元素放置写成带删除的球槽分配问题\cite{aamand2021dynamic,bansal2022deletions}，再用 k-kick、local query router、mini-array 和商压缩拼出一条可以调节的权衡曲线，本文的基础原型主要就是沿这一方向展开的。

\subsection{工程实现与国内研究}

打包压缩哈希表\cite{kurpicz2023pachash}、IcebergHT\cite{pandey2023iceberght,bender2024iceberg}等方案会在真实机器上测实现效果，除了吞吐和装填率，还要把缓存怎么用、重构要花多少代价这些工程因素算进去；国内也有相关工作，比如持久化键值库上的哈希索引优化\cite{wang2024hotaware}、GPU 上的动态哈希表索引\cite{xiong2022starfish}，分别从热点怎么分布、键摘要怎么压、显存怎么分页管理等方向，试着把索引占用的空间降下来。

\section{本文主要工作}

本文在 Bender 等人的论文\cite{bender2022otsh}理论框架下实现 facility--cubby--slot 分层结构的简化原型，完整理论结构中的部分细节则按原型实验需要进行了简化。facility 负责表级分区，cubby 负责局部存储和局部元数据存储，slot 存放商压缩后的高位数据。mini-array 用于维护变长元数据；cubby 内的 k-kick 完成有界交换插入；local query router 将查询从探测序列扫描转化为局部映射访问；mini-array \(B\) 保存键恢复所需的差分信息。该原型保留地址拆分、有限踢出、局部路由和键恢复四个环节。通过实验，可以分别测量查询路径、元素搬移和层级重构带来的开销。

原始多层级 cubby 结构在落地时需要维护层级合并、拆分和删除回填。在完成基础原型后，本文进一步补充冷热分层键存储方案：tail cubby 作为近期数据入口，写满后把较冷的键批量下沉到 SQLite3，这里较冷的键指的是其访问频次相对较低。这个改造借鉴了日志结构合并树先写内存、再批量下沉的思路\cite{oneil1996lsm,dong2021rocksdb}，但本文实现不再完整地保留原结构中多层级 cubby 结构，只保留 tail cubby。第四章用于检验多层级原型结构是否能按规则运行，第五章则单独讨论冷热分层改造与两个对照结构之间的差异。

围绕上述实现，实验部分关注四个问题：(1) 理论模块如何落到具体数据结构，并在插入、查询、删除过程中保持状态一致；(2) mini-array、k-kick 与 local query router 实现并组成完整的哈希系统后，插入、查询与删除各阶段的延迟处于何种量级；(3) 在不同的 \(n\)、\(K\)、\(k\) 分组实验下，k-kick 踢出次数、cubby 拆分与合并次数等指标如何变化，单次搬移次数是否仍受 \(k\) 控制；(4) 冷热分层改造后，内存 tail cubby 与 SQLite3 冷数据层在写入、查询、删除和空间占用效果如何。

\section{本文组织结构}

本章介绍实验背景及意义、国内外哈希表研究的现状以及本文的工作。第\ref{ch:theo_base}章给出相关理论基础以及冷热分层存储思想；第\ref{ch:system}章说明 facility、mini-array、cubby、k-kick、local query router 与键恢复过程的实现细节和冷热分层存储的设计方案。第\ref{ch:exp}章测量多层级原型结构在不同参数下的延迟、踢出和层级重构行为；第\ref{ch:hotcold}章讨论冷热分层改造在写后读和近期访问情况下，插入、删除、查询和空间占用的表现。第\ref{ch:concl}章总结本文工作，并阐述未来优化方向。
